The design of polymer electrolytes requires navigating a coupled chemical and transport landscape in which the most useful candidates may lie close to, or beyond, the best observed conductivity. Generative models promise to reduce this search burden by proposing polymer repeat units conditioned on a target property. Yet a model can appear controllable while merely reweighting structures already represented in its training distribution. A convincing claim of extrapolation therefore requires an explicit boundary, an evaluation region outside that boundary and a verifier that is not the generator objective itself.

We revisit this distinction in a reproducible PolyGen benchmark. The study is deliberately narrow: a character-level autoregressive generator is trained from scratch on a single high-throughput molecular-dynamics (HTP-MD) source, and generated structures are evaluated with a frozen random-forest (RF) surrogate. This design does not establish experimental discovery. Instead, it isolates what conditioning can and cannot do under a fixed data and representation regime.

Our contribution is a four-stage evaluation ladder: interpolation within the training support, upper-tail control near the observed boundary, true numerical extrapolation beyond the training maximum, and structural out-of-distribution evaluation. The first three stages are completed here; the last remains blocked by missing scaffold and family annotations. The central result is a clean separation between controllability and extrapolation. Continuous conditioning improves in-support target adherence, but both continuous and binary conditioning fail the strict frontier-crossing test at every extrapolation query. A six-property joint-conditioning experiment, motivated by recent multi-property molecular generation, reproduces the same limitation.

This negative result is scientifically useful. It identifies the current bottleneck as a limitation of data coverage and representation, rather than a failure to tune the conditioning label finely enough. It also gives a concrete standard for future work: claims of property extrapolation must report support, validity, uncertainty, and an independently defined frontier-crossing criterion.
